\chapter{MLOps und der Modell-Lebenszyklus}
\label{chap:mlops_modelllebenszyklus}

% ============================================================
% AUTOR-NOTIZ STATT LERNZIELE
% ============================================================
\autornotiz{
2024: Production Incident bei einem Support-Bot. Prompt-Update deployed, Golden Set grün. Aber in Produktion: Halluzinationsrate 8\% -> 23\% innerhalb 2h.
Ursache: Provider hatte stillschweigend gpt-4o-mini auf neue Version gehoben (2024-07-18). Unser Eval läuft gegen festes Modell, nicht gegen "latest".
Fix: Model Pinning (gpt-4o-mini-2024-07-18) + Canary Deploy (5\% Traffic) + Auto-Rollback bei Drift > 5\%.
Lektion: Ohne Model Pinning + Canary + Drift Alert fliegst du blind. "Latest" ist kein Deployment-Strategy.
}

% ============================================================
% MOTIVATION
% ============================================================
\section{Motivation}

LLM-Anwendungen bestehen aus mehr als Code. Sie kombinieren Modelle, Daten, Prompts und Infrastruktur. Klassisches Software-Engineering reicht nicht aus: Es braucht einen Lebenszyklus, der alle Komponenten versioniert, testet, deployed und überwacht.

MLOps operationalisiert diese Kette. Für AI Engineers ist sie entscheidend, weil sie aus einem Prototypen eine stabile, skalierbare Produktionsanwendung macht.

% ============================================================
% GRUNDLAGEN
% ============================================================
\section{Was macht MLOps anders?}

Im Vergleich zur klassischen Softwareentwicklung sind drei Aspekte zentral:

\begin{enumerate}[leftmargin=*]
    \item \textbf{Nicht-deterministisches Verhalten:} Der gleiche Prompt liefert nicht immer denselben Output. Tests müssen probabilistisch und robust gestaltet werden.
    \item \textbf{Externe Modellabhängigkeit:} Viele Anwendungen nutzen Managed APIs. Verfügbarkeit, Latenz und Kosten liegen oft außerhalb der eigenen Kontrolle.
    \item \textbf{Daten- und Modelldrift:} Das reale Nutzerverhalten ändert sich, und promptbasierter Code altert schneller als klassischer Application Code.
\end{enumerate}

% ============================================================
% ARCHITEKTUR DES LEBENSZYKLUS
% ============================================================
\section{Der LLMOps-Lebenszyklus}

Eine LLMOps-Architektur besteht aus fünf Hauptphasen:

\begin{itemize}

    \item \textbf{Entwickeln} Prompts, Modell-Konfigurationen, Daten-Pipelines und Service-Code erstellen
    \item \textbf{Testen} Golden Datasets, Prompt-Tests, Kosten- und Latenzprüfungen ausführen
    \item \textbf{Ausrollen} Versionierte Releases, Canary-Deployments und Rollbacks verwalten
    \item \textbf{Überwachen} Drift, Feedback, Kosten und Performance in Produktion beobachten
    \item \textbf{Verbessern} Regeln, Daten und Prompts iterativ anpassen und den Kreislauf erneuern
\end{itemize}

\subsection{Systemkomponenten}

\begin{itemize}[leftmargin=*]
    \item \textbf{Versionierung}: Git für Code und Prompts, DVC/LakeFS für Trainingsdaten, Model Cards für Modelle
    \item \textbf{CI/CD}: Automatisierte Tests, Review-Gates, Canary-Deployments
    \item \textbf{Observability}: Metriken, Tracing, Log-Events, Drift-Metriken
    \item \textbf{Governance}: Verantwortlichkeiten, Dokumentation, Audit-Logs
    \item \textbf{Feedback}: Nutzerfeedback, menschliche Evaluierung, Incident Reports
\end{itemize}

% ============================================================
% VERSIONIERUNG
% ============================================================
\section{Versionierung von Prompts, Daten und Modellen}

\subsection{Prompts als Code}

Prompts gehören in die Codebasis. Jede Änderung braucht eine Git-Historie, ein Review und muss zurücksetzbar sein:

\begin{itemize}[leftmargin=*]
    \item \textbf{Prompt-Templates} als parametrische Dateien
    \item \textbf{Prompt-Tests} gegen Golden Datasets
    \item \textbf{Release Notes} für Prompt-Änderungen
\end{itemize}

\praxishinweis{Prompts in der Codebasis zu versionieren reicht nicht -- du brauchst auch die zugehörigen Test-Ergebnisse. Ein Prompt-Change kann dieselben Tests bestehen, aber Produktionsoutput-Qualität ändern. Speichere Evaluations-Results immer mit der Prompt-Version.}

\subsection{Datenversionierung}

Daten für Evaluationen, Referenzen und semantische Indizes müssen versioniert werden:

\begin{itemize}[leftmargin=*]
    \item \textbf{Daten-Snapshots} für historische Vergleichbarkeit
    \item \textbf{Schema-Checks} zur Erkennung von Drift
    \item \textbf{Reproduzierbare Builds} durch festgehaltene Datenversionen
\end{itemize}

\subsection{Modelldokumentation}

Jedes Modell braucht eine \textbf{Model Card}:

\begin{itemize}[leftmargin=*]
    \item Zweck und Einsatzgebiet
    \item Trainingsdaten- oder Prompt-Basis
    \item bekannte Schwächen und Risiken
    \item Performance-Metriken und Kosten
\end{itemize}

% ============================================================
% CICD
% ============================================================
\section{CI/CD für LLM-Systeme}

Eine LLMOps-Pipeline enthält neben Code-Checks auch prompt- und datenorientierte Schritte.

\subsection{Beispiel-Pipeline}

\begin{itemize}[leftmargin=*]
    \item \textbf{Linting}\ für Code, Prompts und JSON-Schema
    \item \textbf{Golden-Dataset-Tests}\ zur Qualitätsabsicherung
    \item \textbf{LLM-Judge-Tests}\ zur semantischen Bewertung
    \item \textbf{Kosten-/Latenz-Checks}\ zur Budgetkontrolle
    \item \textbf{Canary Deploy}\ zur risikoreduzierten Einführung
\end{itemize}

\subsection{Praktischer Ablauf}

Eine CI/CD-Stage kann so aussehen:

\begin{lstlisting}[language=Bash,caption={Beispiel-Workflow für LLMOps in GitHub Actions}]
name: LLMOps Pipeline
on: [pull_request, push]

jobs:
  test:
    runs-on: ubuntu-latest
    steps:
      - uses: actions/checkout@v4
      - name: Install dependencies
        run: pip install -r requirements.txt
      - name: Prompt lint
        run: python scripts/prompt_lint.py prompts/
      - name: Run golden dataset tests
        run: python scripts/evaluate_golden.py
      - name: Validate JSON schema
        run: python scripts/validate_output_schema.py

  canary:
    needs: test
    if: github.event_name == 'push' && github.ref == 'refs/heads/main'
    runs-on: ubuntu-latest
    steps:
      - uses: actions/checkout@v4
      - name: Deploy canary
        run: ./scripts/deploy_canary.sh
      - name: Run canary smoke tests
        run: python scripts/canary_smoke_test.py
\end{lstlisting}

\praxishinweis{CI/CD für LLMs braucht Cost-Gates. Definiere einen maximalen Token-Verbrauch pro Test-Suite. Überschreitet ein Prompt-Update diesen Wert, schlägt die Pipeline fehl -- bevor die Kosten in Produktion explodieren.}

% ============================================================
% MONITORING
% ============================================================
\section{Monitoring und Drift-Erkennung}

Überwachung bedeutet nicht nur, dass ein Dienst erreichbar ist. Für LLM-Systeme zählen auch Inhalt und Qualität.

\subsection{Wichtige Produktionsmetriken}

\begin{itemize}

    \item \textbf{Qualitäts-Drift} Veränderung der Antwortqualität über Zeit. Gemessen durch Stichproben, Feedback oder automatisierte Judges.
    \item \textbf{Daten-Drift} Veränderungen in den Eingabedaten (Länge, Domäne, Token-Verteilung).
    \item \textbf{Kosten-Drift} Unerwartete Steigerungen durch höhere Tokens, längere Ausgaben oder neue Nutzerprofile.
    \item \textbf{Performance-Drift} Steigende Latenzen, höhere Fehlerquoten oder vermehrte Timeouts.
\end{itemize}

\subsection{Alerting}

Alerts sollten nicht nur bei Ausfällen ausgelöst werden, sondern auch bei:

\begin{itemize}[leftmargin=*]
    \item \textbf{Signifikanter Kostenabweichung} im Vergleich zur Baseline
    \item \textbf{Spürbarer Qualitätsverschlechterung} im Golden Dataset
    \item \textbf{Datenanomalien} bei Eingabeverteilungen
    \item \textbf{Anstieg von User Complaints} oder Ablehnungen
\end{itemize}

% ============================================================
% GOVERNANCE
% ============================================================
\section{Governance und Verantwortlichkeiten}

Eine stabile KI-Produktion braucht klare Rollen:

\begin{itemize}[leftmargin=*]
    \item \textbf{Model Owner}: Verantwortlich für die Modellwahl, Release-Planung und Dokumentation
    \item \textbf{Data Steward}: Sorgt für Datenqualität, Versionierung und Compliance
    \item \textbf{Product Owner}: Bewertet Nutzerfeedback und priorisiert Verbesserungen
    \item \textbf{Security Owner}: Überprüft Zugriff, Datenschutz und Missbrauchsschutz
\end{itemize}

\subsection{Audit und Dokumentation}

Jede Modelländerung sollte nachvollziehbar sein:

\begin{itemize}[leftmargin=*]
    \item Release Notes für Prompts, Modelle und Indexdaten
    \item Logs von Evaluationen und Governance Reviews
    \item Model Cards und Risikoanalysen
    \item Approval-Gates vor Deployments
\end{itemize}

% ============================================================
% PROMPT-VERSIONIERUNG
% ============================================================
\section{Prompt-Versionierung in der Praxis}

\subsection{Strukturierte Prompt-Ordner}

Eine bewährte Struktur trennt Prompts nach Umgebung und Domäne:

\begin{verbatim}
prompts/
  staging/
    customer-support/
      v1__de-escalation.jinja2
      v2__de-escalation.jinja2
    code-assistant/
      v3__code-review.jinja2
  production/
    customer-support/
      v2__de-escalation.jinja2
    code-assistant/
      v3__code-review.jinja2
\end{verbatim}

Jeder Prompt ist eine Jinja2-Vorlage in Git. Ein Wechsel der aktiven Version erfolgt durch Deployment -- nicht durch Datei-Edit.

\subsection{Automatisierte Prompt-Tests}

Vor dem Deployment durchläuft jeder Prompt eine automatisierte Testsuite:

\begin{lstlisting}[language=Python,caption=Prompt-Test mit Golden Dataset]
import pytest
from jinja2 import Template

GOLDEN = [
    {"input": "Wie kuendige ich mein Abo?",
     "expected_keywords": ["kuendigen", "Abo", "Optionen"]},
    {"input": "Ich habe meine Rechnung nicht erhalten",
     "expected_keywords": ["Rechnung", "erneut", "Support"]},
]

def test_prompt_golden():
    template = Template(open("prompts/support.jinja2").read())
    for case in GOLDEN:
        result = template.render(query=case["input"])
        for kw in case["expected_keywords"]:
            assert kw in result, \
                f"Fehlendes Schluesselwort {kw} in Prompt-Output"
\end{lstlisting}

\subsection{Rollback-Strategie}

Jeder Prompt-Deployment erzeugt einen versionierten Snapshot. Ein Rollback bedeutet: vorherigen Snapshot als aktiv markieren und erneut deployen. Wichtig: die zugehörigen Embeddings und Konfigurationen müssen ebenfalls zurückgesetzt werden können.

% ============================================================
% KOSTENMANAGEMENT
% ============================================================
\section{Kostenmanagement und Budgetkontrolle}

LLM-Kosten sind kein Afterthought -- sie können ein Produkt unwirtschaftlich machen.

\subsection{Kostentreiber identifizieren}

Die häufigsten Kostentreiber:

\begin{itemize}

    \item \textbf{Kontextfenster} Längere Prompts bedeuten höhere Kosten pro Anfrage. Ein System-Prompt mit 2000 Tokens kostet bei jeder Interaktion -- unabhängig vom Nutzer-Prompt.
    \item \textbf{Output-Länge} Modelle, die ausführliche Antworten generieren, verbrauchen mehr Output-Tokens. Besonders kritisch bei Agenten mit mehreren Tool-Aufrufen.
    \item \textbf{Fehlertoleranz} Wiederholte Anfragen durch Timeouts oder fehlerhafte Validierung vervielfachen die Kosten.
    \item \textbf{Modellwahl} Teurere Modelle pro Token sind nicht immer nötig -- ein Semantic Router kann einfache Anfragen an günstigere Modelle delegieren.
\end{itemize}

\subsection{Kostenbudgetierung}

Setze ein mehrstufiges Budgetsystem auf:

\begin{enumerate}[leftmargin=*]
    \item \textbf{Harter Deckel} (Hard Cap): Maximale monatliche Ausgaben. Bei Überschreitung wird auf ein günstigeres Modell oder Caching-only-Modus umgeschaltet.
    \item \textbf{Weiche Warnung} (Soft Alert): Bei 80\,\% des Budgets wird ein Alarm ausgelöst.
    \item \textbf{Pro-Kunde-Limit}: Budgetbegrenzung pro Mandant oder API-Key, um Kostenexplosion durch einzelne Nutzer zu verhindern.
\end{enumerate}

\begin{lstlisting}[language=Python,caption=Kosten-Tracker für LLM-API-Aufrufe]
import time
from dataclasses import dataclass, field

@dataclass
class CostTracker:
    budget_limit: float
    current_cost: float = 0.0
    alerts: list = field(default_factory=list)

    def track(self, model: str, prompt_tokens: int,
              completion_tokens: int, prices: dict) -> None:
        price = prices.get(model, 0)
        cost = (prompt_tokens * price["input"]
                + completion_tokens * price["output"]) / 1_000_000
        self.current_cost += cost
        ratio = self.current_cost / self.budget_limit
        if ratio >= 0.8:
            self.alerts.append(
                f"Warnung: {ratio:.0%} des Budgets erreicht"
            )

    def over_budget(self) -> bool:
        return self.current_cost >= self.budget_limit
\end{lstlisting}

\subsection{Kostenoptimierende Pattern}

\begin{itemize}[leftmargin=*]
    \item \textbf{Prompt Caching}: Wiederkehrende Prefix-Tokens werden gecached und nur einmal berechnet.
    \item \textbf{Batch-Processing}: Asynchrone Anfragen zusammenfassen reduziert API-Calls.
    \item \textbf{Output-Limits}: \texttt{max\_tokens} strikt setzen -- auch wenn das Modell weiterreden könnte.
    \item \textbf{Caching auf Applikationsebene}: Semantisch ähnliche Anfragen können ohne API-Aufruf beantwortet werden (siehe Kapitel 19).
\end{itemize}

\praxishinweis{Setze Kosten-Alarme auf zwei Ebenen: Tagesbudget (hartes Limit) und Run-Rate (gleitender 7-Tage-Durchschnitt). Ein Tagesspitze ist ok, eine steigende Run-Rate zeigt ein strukturelles Problem.}

% ============================================================
% INCIDENT RESPONSE
% ============================================================
\section{Incident Response und Rollback}

LLM-Produktionssysteme fallen anders aus als klassische APIs.

\subsection{Typische Incidents}

\begin{itemize}

    \item \textbf{Qualitätseinbruch} Das Modell liefert plötzlich schlechtere Antworten -- durch ein Modell-Update des Providers oder Drift in den Eingabedaten.
    \item \textbf{Kosten-Explosion} Ein fehlerhafter Prompt oder eine Endlosschleife in einem Agenten verursacht unerwartet hohe Token-Verbräuche.
    \item \textbf{Security Incident} Prompt Injection, Jailbreaking oder Datenexfiltration über die Ausgabe.
    \item \textbf{Provider-Ausfall} Der LLM-Provider ist nicht erreichbar oder liefert konsistent Fehler zurück.
\end{itemize}

\subsection{Runbook-Struktur}

Jeder Incident-Typ braucht ein Runbook:

\begin{enumerate}[leftmargin=*]
    \item \textbf{Detect}: Metriken/Alerts lösen aus (Kostenanstieg, Fehlerrate, Qualitätsscore)
    \item \textbf{Triage}: Ist es ein Provider- oder Systemproblem? Betrifft es alle Nutzer oder eine Gruppe?
    \item \textbf{Mitigate}: Fallback-Modell aktivieren, Cache leeren, Prompt auf letzte stabile Version zurücksetzen
    \item \textbf{Resolve}: Root-Cause-Analyse, Fix deployen, Monitoring verstärken
    \item \textbf{Learn}: Runbook aktualisieren, Metriken nachschärfen, Post-Mortem schreiben
\end{enumerate}

\subsection{Fallback-Strategien}

\begin{itemize}[leftmargin=*]
    \item \textbf{Multi-Provider}: Bei Ausfall eines Providers auf einen zweiten umschalten (z.\,B. OpenAI $\rightarrow$ Anthropic)
    \item \textbf{Modell-Degradation}: Bei Überlastung auf ein kleineres, günstigeres Modell wechseln
    \item \textbf{Cache-Only-Modus}: Bei komplettem Provider-Ausfall nur noch gecachte Antworten ausliefern -- mit deutlichem Hinweis für den Nutzer
\end{itemize}

% ============================================================
% ZUSAMMENFASSUNG
% ============================================================
\begin{zusammenfassungEnv}
MLOps macht LLM-Anwendungen produktionsreif: nicht nur durch Deployment, sondern durch Versionierung, Tests, Monitoring, Kostenkontrolle und Incident Response. Ein strukturierter Lebenszyklus reduziert Risiko, Kosten und Drift, und ermöglicht wiederholbare Verbesserungen.
\end{zusammenfassungEnv}




\par\mbox{}\par
\begin{merkeEnv}
LLM-Deployment ist nicht wie klassisches Web-Deployment. External Dependencies, probabilistisches Verhalten und Kosten pro Request erfordern neue Patterns: Circuit-Breaker, Caching, Multi-Provider-Fallback und kontinuierliches Cost-Monitoring.
\end{merkeEnv}

% ============================================================
% PRAXISPROJEKT
% ============================================================
\section{Praxisprojekt -- MLOps-Pipeline für einen LLM-Chat-Service}

\textbf{Ziel:} Baue eine vollständige MLOps-Pipeline für einen Chat-Service.

\textbf{Aufgaben:}
\begin{enumerate}[leftmargin=*]
    \item Versioniere Prompts, Golden Dataset und Evaluation-Scripts in Git
    \item Konfiguriere GitHub Actions: Linting $\rightarrow$ Golden Dataset Tests $\rightarrow$ Cost Check $\rightarrow$ Canary Deploy
    \item Implementiere Prometheus-Metriken: Requests, Latenz, Kosten, Error Rate
    \item Richte Grafana-Dashboard ein: Quality Score, Cost/Query, P95 Latency
    \item Definiere SLAs und Alerts: Quality $<$ 0.7, Cost/Query $>$ 0.02 EUR, P95 $>$ 3s
    \item Schreibe Runbook für drei Incident-Typen: Quality Drop, Cost Spike, Provider Outage
    \item Führe einen Fire Drill durch: Simuliere Provider-Ausfall, miss Fallback, messe Recovery-Time
\end{enumerate}

\textbf{Erfolgskriterium:} Pipeline läuft grün bei jedem Push. Canary-Deployment bei Merge auf main. Alerts feuern bei injizierten Anomalien innerhalb von 2 Minuten.

% ============================================================
% RESSOURCEN
% ============================================================
\section{Ressourcen}

\begin{itemize}[leftmargin=*]
    \item \textbf{MLflow}: \href{https://mlflow.org}{mlflow.org} -- Open-Source ML Lifecycle Management
    \item \textbf{Langfuse}: \href{https://langfuse.com}{langfuse.com} -- LLM Observability \& Evaluation
    \item \textbf{Evidently AI}: \href{https://www.evidentlyai.com}{evidentlyai.com} -- ML-Monitoring und Drift-Erkennung
\end{itemize}

\textbf{Nächster Schritt:} Observability läuft -- jetzt Security: Input/Output Guards, PII-Masking, Injection Defense, Agent Tool-Whitelist, HITL für kritische Aktionen, Incident Runbooks. SupportPilot's Safety-Layer.
\endinput

\clearpage
